\documentclass[a4paper,11pt]{article}

%%% Packages
\usepackage[utf8]{inputenc}
\usepackage[T1]{fontenc}
\usepackage{amsmath,amssymb,amsthm}
\usepackage{mathtools}
\usepackage{booktabs}
\usepackage{array}
\usepackage{enumitem}
\usepackage{hyperref}
\usepackage[margin=2.5cm]{geometry}
\usepackage{xcolor}

%%% Theorem environments
\newtheorem{theorem}{Theorem}[section]
\newtheorem{lemma}[theorem]{Lemma}
\newtheorem{proposition}[theorem]{Proposition}
\newtheorem{corollary}[theorem]{Corollary}
\newtheorem{claim}[theorem]{Claim}
\newtheorem{conjecture}[theorem]{Conjecture}
\theoremstyle{definition}
\newtheorem{definition}[theorem]{Definition}
\newtheorem{example}[theorem]{Example}
\newtheorem{remark}[theorem]{Remark}

%%% Notation macros
\newcommand{\ALC}{\ensuremath{\mathcal{ALC}}}
\newcommand{\ALCI}{\ensuremath{\mathcal{ALCI}}}
\newcommand{\ALCIRCC}[1]{\ensuremath{\mathcal{ALCI}_{\mathit{RCC#1}}}}
\newcommand{\NR}{\ensuremath{N_R}}
\newcommand{\NC}{\ensuremath{N_C}}
\newcommand{\DeltaI}{\ensuremath{\Delta^{\mathcal{I}}}}
\newcommand{\interp}{\ensuremath{\mathcal{I}}}
\newcommand{\DR}{\ensuremath{\mathsf{DR}}}
\newcommand{\PO}{\ensuremath{\mathsf{PO}}}
\newcommand{\EQ}{\ensuremath{\mathsf{EQ}}}
\newcommand{\PP}{\ensuremath{\mathsf{PP}}}
\newcommand{\PPI}{\ensuremath{\mathsf{PPI}}}
\newcommand{\inv}{\ensuremath{\mathsf{inv}}}
\newcommand{\comp}{\ensuremath{\mathsf{comp}}}
\newcommand{\cl}{\ensuremath{\mathsf{cl}}}
\newcommand{\tp}{\ensuremath{\mathsf{tp}}}
\newcommand{\EXPTIME}{\textsc{ExpTime}}


\title{Direct Soundness for $\ALCIRCC{5}$:\\
  Demand Satisfaction without the Henkin Construction}

\author{Claude\thanks{Claude is an AI assistant by Anthropic.
  This research was prompted by Michael Wessel (\texttt{miacwess@gmail.com}),
  who introduced the $\ALCIRCC{}$ family in his doctoral work at the
  University of Hamburg.}}

\date{April 2026}

\begin{document}
\maketitle

\begin{center}
\fbox{\parbox{0.92\textwidth}{%
\textbf{Disclaimer.}\; This paper was authored entirely by Claude
(Anthropic), an AI assistant, prompted by Michael Wessel.  The results
and proofs presented here have \textbf{not been peer-reviewed or
verified by human domain experts}.  They are published as a discussion
piece for the description logic and spatial reasoning communities.}}
\end{center}

\medskip

\begin{center}
\fbox{\parbox{0.92\textwidth}{%
\textbf{Status (April 2026).}\;
This paper presents a new approach to the soundness problem for the
$\ALCIRCC{5}$ tableau, bypassing the Henkin construction and its
extension gap.  The approach is supported by exhaustive computational
verification but the formal proof is \textbf{not complete}---the
cross-chain edge assignment (Section~\ref{sec:cross-edges}) remains
a conjecture, equivalent to the Extension Solvability Conjecture
of~\cite{TrianglePaper}.  The novel contribution is the clean
decomposition into finite demands (fully verified) and infinite
chains (structurally rigid), which significantly narrows the remaining
gap.}}
\end{center}


\bigskip

%% ============================================================
\begin{abstract}
The companion paper~\cite{CompanionPaper} presented a tableau calculus
for concept satisfiability in $\ALCIRCC{5}$ with established
termination and completeness, but left soundness open due to the
\emph{extension gap}: the Henkin construction's disjunctive CSP is
not necessarily path-consistent.

We propose a direct soundness approach that bypasses the Henkin
construction entirely.  Using exhaustive computational analysis of
the RCC5 composition table, we establish a clean trichotomy for
existential demand satisfaction:
\begin{enumerate}[nosep]
  \item \textbf{No PP/PPI demands}: the completion graph directly
    satisfies all demands (zero failures in all tested configurations).
  \item \textbf{Acyclic PP/PPI chains}: demands resolve after
    bounded extra expansion (depth $\leq$ number of Hintikka types).
  \item \textbf{PP/PPI cycles}: require infinite chains, but these
    are algebraically rigid ($\comp(\PP,\PP) = \{\PP\}$,
    $\comp(\PPI,\PPI) = \{\PPI\}$) and structurally trivial.
\end{enumerate}
The root cause of all demand failures is a single algebraic fact:
$\PPI \notin \comp(\DR, \PPI) = \{\DR\}$, the unique self-absorption
failure in the RCC5 composition table.

Combined with the tableau's termination and completeness
from~\cite{CompanionPaper}, this reduces the decidability of
$\ALCIRCC{5}$ to a single remaining question: the consistent
assignment of cross-chain edges, which is precisely the Extension
Solvability Conjecture of~\cite{TrianglePaper}.
\end{abstract}

\medskip
\noindent\textbf{Keywords:} description logic, spatial reasoning,
Region Connection Calculus, decidability, demand satisfaction,
self-absorption


%% ============================================================
\section{Introduction}
\label{sec:intro}

\subsection{The extension gap}

The decidability of concept satisfiability in $\ALCIRCC{5}$ has
resisted proof despite multiple approaches.  The companion
paper~\cite{CompanionPaper} established termination and completeness
of a tableau calculus with equality anywhere-blocking, but
soundness---the direction from an open completion graph to a
model---remained open.

All previous soundness attempts relied on the \emph{Henkin
construction}: extract a quasimodel from the open completion graph,
unravel it into an infinite tree, then assign cross-subtree edges
via a disjunctive RCC5 constraint network.  The \emph{extension gap}
is the observation that this constraint network is not necessarily
path-consistent, so the full tractability theorem for RCC5 does not
directly apply.

\subsection{The direct approach}

We propose bypassing the Henkin construction entirely.  Instead of
unraveling and re-assigning edges, we ask: does the open completion
graph \emph{already satisfy} all existential demands?

An expanded (non-blocked) node trivially has its demands satisfied
by its own successors.  The question is about \emph{blocked nodes}:
a node $b$ blocked by $a$ (with $\mathcal{L}(b) = \mathcal{L}(a)$)
has the same existential demands as $a$, but was not expanded.
Are $b$'s demands satisfied by other nodes already in the graph?

If yes, the completion graph---or a bounded expansion of it---is
directly a model, with no Henkin step.

\subsection{Main findings}

Through exhaustive computational analysis, we establish:

\begin{enumerate}[nosep]
  \item \textbf{Unique algebraic obstruction} (Section~\ref{sec:self-absorption}).
    The RCC5 composition table has exactly one self-absorption failure:
    $\PPI \notin \comp(\DR, \PPI) = \{\DR\}$.  This means: if a
    blocked node $b$ has relation $\DR$ to its blocker $a$, then
    $b$ cannot reach $a$'s $\PPI$-successors via $\PPI$.  All other
    demand/relation combinations work.

  \item \textbf{Rigid chain classification} (Section~\ref{sec:rigid}).
    $\PP$ and $\PPI$ are the only relations with
    $\comp(R,R) = \{R\}$ (``rigid chains'').  $\DR$ and $\PO$
    satisfy $\comp(R,R) = \{\DR,\PO,\PP,\PPI\}$ (``diluting''),
    meaning that after one hop, all four relations become possible
    and demands can be redirected.

  \item \textbf{Trichotomy of demand satisfaction}
    (Section~\ref{sec:trichotomy}).  All failures decompose into
    three cases based on the PP/PPI structure of the demand graph.
\end{enumerate}


%% ============================================================
\section{Preliminaries}
\label{sec:prelim}

We recall the essential definitions.  The logic $\ALCIRCC{5}$ has
role names $\NR = \{\DR, \PO, \EQ, \PP, \PPI\}$ with the standard
RCC5 composition table.  Interpretations are complete graphs:
every pair of distinct elements is related by exactly one base
relation from $\NR \setminus \{\EQ\}$ (strong EQ semantics).

After EQ-normalization, we work with the four non-EQ relations
$\NR^{-} = \{\DR, \PO, \PP, \PPI\}$.  The \emph{closure}
$\cl(C_0)$ of a concept $C_0$ is finite and computable.  A
\emph{Hintikka type} $\tau$ is a maximally consistent subset of
$\cl(C_0)$.  The \emph{existential demands} of $\tau$ are
$\mathsf{dem}(\tau) = \{(R, \tau') : \exists R.D \in \tau,\;
D \in \tau'\}$ for some pairing of demands to target types.

\begin{definition}[Demand satisfaction]
A node $v$ in a labeled complete RCC5 graph $\mathcal{G}$ has
\emph{demand $(R, \tau')$ satisfied} if there exists $w \neq v$
with $\mathcal{L}(w) = \tau'$ and $\mathcal{E}(v, w) = R$.

A graph has \emph{uniform demand satisfaction} for type $\tau$
and demand $(R, \tau')$ if every node of type $\tau$ has $(R,\tau')$
satisfied.
\end{definition}

\begin{definition}[PP/PPI demand graph]
Given a set of types with demands, the \emph{PP/PPI demand graph}
has types as vertices and an edge $\tau \to \tau'$ whenever
$(\PP, \tau') \in \mathsf{dem}(\tau)$ or
$(\PPI, \tau') \in \mathsf{dem}(\tau)$.
\end{definition}


%% ============================================================
\section{The Self-Absorption Analysis}
\label{sec:self-absorption}

\begin{definition}[Self-absorption]
A base relation $R \in \NR^{-}$ is \emph{self-absorbing under}
$S \in \NR^{-}$ if $R \in \comp(S, R)$.
\end{definition}

Self-absorption captures the key property for demand satisfaction:
if two nodes of the same type have relation $S$ between them, and
one has an $R$-successor $w$, can the other also have $R$ to $w$?
This requires $R \in \comp(S, R)$.

\begin{proposition}[Unique failure]
\label{prop:unique-failure}
The RCC5 composition table has exactly one self-absorption failure:
\[
  \PPI \notin \comp(\DR, \PPI) = \{\DR\}.
\]
All other 15 pairs $(R, S) \in \NR^{-} \times \NR^{-}$ satisfy
$R \in \comp(S, R)$.
\end{proposition}

\begin{proof}
Exhaustive verification against the RCC5 composition table.  The
full 16-entry check:

\medskip
\begin{center}
\renewcommand{\arraystretch}{1.1}
\begin{tabular}{c|cccc}
$R \in \comp(S,R)$? & $S{=}\DR$ & $S{=}\PO$ & $S{=}\PP$ & $S{=}\PPI$ \\
\hline
$R{=}\DR$  & \checkmark & \checkmark & \checkmark & \checkmark \\
$R{=}\PO$  & \checkmark & \checkmark & \checkmark & \checkmark \\
$R{=}\PP$  & \checkmark & \checkmark & \checkmark & \checkmark \\
$R{=}\PPI$ & $\times$   & \checkmark & \checkmark & \checkmark \\
\end{tabular}
\end{center}
\medskip

The single failure: $\comp(\DR, \PPI) = \{\DR\}$, so
$\PPI \notin \comp(\DR, \PPI)$. \qedhere
\end{proof}

\begin{remark}
The dual failure $\PP \notin \comp(\PP, \DR) = \{\DR\}$ describes
the same algebraic obstruction from the inverse perspective (take
inverses throughout).  It is not a separate self-absorption failure,
as self-absorption is defined as $R \in \comp(S, R)$, not
$R \in \comp(R, S)$.
\end{remark}


%% ============================================================
\section{Rigid vs.\ Diluting Relations}
\label{sec:rigid}

\begin{definition}[Rigidity]
A relation $R \in \NR^{-}$ is \emph{rigid} if $\comp(R, R) = \{R\}$.
It is \emph{diluting} if $\comp(R, R) = \NR^{-}$.
\end{definition}

\begin{proposition}
\label{prop:rigid}
$\PP$ and $\PPI$ are rigid.  $\DR$ and $\PO$ are diluting.
\end{proposition}

\begin{proof}
Direct computation:
\begin{align*}
  \comp(\PP, \PP)   &= \{\PP\},  &
  \comp(\PPI, \PPI) &= \{\PPI\}, \\
  \comp(\DR, \DR)   &= \{\DR, \PO, \PP, \PPI\},  &
  \comp(\PO, \PO)   &= \{\DR, \PO, \PP, \PPI\}. \qedhere
\end{align*}
\end{proof}

The rigidity of $\PP$ and $\PPI$ means that chains of these
relations are ``algebraically frozen'': if $\PP(d_0, d_1)$ and
$\PP(d_1, d_2)$, then $\PP(d_0, d_2)$ is forced.  Every element
in a $\PP$-chain has $\PP$ to all later elements and $\PPI$ to
all earlier ones.  No other relation can appear between chain
elements.

In contrast, the dilution of $\DR$ and $\PO$ means that after
one hop, all four relations become possible.  This is why demands
for $\DR$ and $\PO$ can always be redirected through alternative
paths in the graph.

\begin{corollary}
\label{cor:dilution}
If the blocked node $b$ has relation $\DR$ to its blocker $a$,
and $a$ has an $R$-successor $w$ with $R \in \{\DR, \PO, \PP\}$,
then $R \in \comp(\DR, R)$ and $b$ can have $R$ to $w$.

Only the case $R = \PPI$ fails:
$\PPI \notin \comp(\DR, \PPI) = \{\DR\}$.
\end{corollary}


%% ============================================================
\section{The Trichotomy of Demand Satisfaction}
\label{sec:trichotomy}

We now classify all possible demand satisfaction scenarios based
on the structure of the PP/PPI demand graph.

\subsection{Case 1: No PP/PPI demands}

\begin{theorem}[Direct satisfaction]
\label{thm:case1}
If the demand structure contains no $\PP$ or $\PPI$ demands
\textup{(}all demands use $\DR$ or $\PO$\textup{)}, then every
open completion graph satisfies all demands uniformly.
\end{theorem}

\begin{proof}[Proof sketch]
Since $\DR$ and $\PO$ are self-absorbing under all relations
(Proposition~\ref{prop:unique-failure}), a blocked node $b$ with
relation $S$ to its blocker $a$ satisfies: for any demand
$(R, \tau')$ of $b$'s type with $R \in \{\DR, \PO\}$, the
blocker's $R$-successor $w$ (of type $\tau'$) is also reachable
from $b$ via $R$, since $R \in \comp(S, R)$ holds universally.
\end{proof}

\begin{remark}
Computational verification confirms zero failures across all
tested configurations (up to 5 nodes, 3 types) for demand
structures without PP/PPI demands.
\end{remark}

\subsection{Case 2: Acyclic PP/PPI chains}

\begin{definition}[PP/PPI chain length]
The \emph{PP/PPI chain length} $L$ of a demand structure is the
length of the longest acyclic path in the PP/PPI demand graph.
\end{definition}

\begin{theorem}[Bounded expansion]
\label{thm:case2}
If the PP/PPI demand graph is acyclic with chain length $L$, then
expanding the completion graph $L+1$ levels beyond the blocking
point resolves all demand failures.  In particular, $L$ is bounded
by the number of Hintikka types, which is at most $2^{|\cl(C_0)|}$.
\end{theorem}

\begin{proof}[Proof sketch]
The stuck node $b$ has an unsatisfied $\PPI$-demand to some type
$\tau'$, caused by the $\DR$-to-$\PPI$ self-absorption failure.
Expanding $b$ creates a $\PPI$-successor $e$ of type $\tau'$.
Node $e$'s demands may include further $\PPI$-demands, continuing
the chain.  After at most $L$ expansions, the chain terminates at
a type with no PP/PPI demands (acyclicity), and that terminal
node's demands (all $\DR/\PO$) are satisfied by Case~1.

At each level, the new node has $\PP$ to its parent (since it is
a $\PPI$-successor) and $\DR$ to distant ancestors (via
$\comp(\PP, \DR) = \{\DR\}$).  The terminal node's $\DR/\PO$
demands are satisfied because:
\begin{itemize}[nosep]
  \item $\DR$ to seed nodes is forced by repeated composition
    through $\PP$ and $\DR$ (both yield $\DR$).
  \item $\DR$ is self-absorbing under all relations, so the $\DR$
    demand to any type is satisfied by the corresponding seed node.
\end{itemize}
Each intermediate node's $\PPI$-demand is then satisfied by its
own $\PPI$-successor at the next level.
\end{proof}

\begin{remark}
Computational verification: for 2 types with up to 2 demands each,
all 60 stuck scenarios (no PP/PPI cycle) resolve at depth~2.  For
3 types with 1 demand each, 1170 of 1194 resolve at depth~2 and
the remaining 24 (all PP/PPI chains of length~2) resolve at
depth~3.
\end{remark}

\subsection{Case 3: PP/PPI cycles}

\begin{theorem}[Infinite chains]
\label{thm:case3}
If the PP/PPI demand graph contains a cycle, the demands require
infinite $\PP$ or $\PPI$ chains.  These chains are algebraically
rigid: all pairs of elements within a chain have the same relation
\textup{(}$\PP$ or $\PPI$\textup{)}.
\end{theorem}

\begin{proof}
A cycle $\tau_0 \to \tau_1 \to \cdots \to \tau_k = \tau_0$ in
the PP/PPI demand graph means each $\tau_i$ has a $\PP$ or $\PPI$
demand that must be satisfied.  Expanding this cycle produces an
infinite sequence of nodes, each providing the $\PPI$-witness for
the previous one.

By rigidity ($\comp(\PP, \PP) = \{\PP\}$ and
$\comp(\PPI, \PPI) = \{\PPI\}$), composing multiple chain steps
preserves the relation.  Every pair of nodes separated by $n$
chain steps has relation $\PP^n = \PP$ (or $\PPI^n = \PPI$).
The chain is a totally ordered set under $\PP$.

The simplest cycle is self-referential: type $\tau$ has demand
$(\PPI, \tau)$.  This produces an infinite descending $\PPI$-chain
$d_0 \stackrel{\PPI}{\to} d_1 \stackrel{\PPI}{\to} d_2
\stackrel{\PPI}{\to} \cdots$ where each $d_i$ is a $\PPI$-witness
for $d_{i-1}$.
\end{proof}

\begin{remark}[The chain construction is standard]
Infinite $\PP$-chains are well-known in $\ALCIRCC{5}$: the concept
$\exists \PP.\top \sqcap \forall \PP.(\exists \PP.\top)$ requires
an infinite model~\cite{Wessel2003}.  The chain construction itself
is straightforward; the difficulty lies in assigning \emph{cross-chain
edges}.
\end{remark}


%% ============================================================
\section{Model Construction: The Two-Phase Approach}
\label{sec:construction}

We describe the model construction that results from the trichotomy.

\subsection{Phase 1: Finite skeleton}

Given an open completion graph $\mathcal{G}$ for $C_0$:

\begin{enumerate}[nosep]
  \item Identify the PP/PPI demand graph from the Hintikka types
    in $\mathcal{G}$.
  \item Compute the chain length $L$ (longest acyclic PP/PPI path).
  \item Expand the completion graph $L + 1$ additional levels
    beyond each blocking point, using the standard $\exists$-rule
    and CF-rule.
  \item The resulting expanded graph $\mathcal{G}^+$ satisfies
    all \emph{non-cyclic} demands uniformly (by Theorems~\ref{thm:case1}
    and~\ref{thm:case2}).
\end{enumerate}

The finite skeleton consists of the active nodes of $\mathcal{G}^+$,
with all edges assigned.

\subsection{Phase 2: Infinite chain extensions}

For each PP/PPI cycle in the demand graph:

\begin{enumerate}[nosep]
  \item Identify the cycle types $\tau_0, \tau_1, \ldots, \tau_k$
    and their PP/PPI demands.
  \item For each node $v$ in the finite skeleton whose type
    participates in the cycle, attach an infinite chain:
    $v = d_0 \stackrel{\PPI}{\to} d_1 \stackrel{\PPI}{\to}
    d_2 \stackrel{\PPI}{\to} \cdots$
    with $\tp(d_i) = \tau_{i \bmod (k+1)}$.
  \item Each chain element $d_i$ has $\PPI$ to $d_{i+1}$
    (satisfying the cyclic demand) and $\PP$ to $d_{i-1}, \ldots,
    d_0$ (by rigidity of $\PPI$).
\end{enumerate}

Within each chain, all edges are determined by rigidity.  The
chain elements' non-PP/PPI demands are satisfied by the finite
skeleton (their DR/PO demands reach the skeleton via composition
through PP and DR, which yields DR, and DR is self-absorbing).


\subsection{Phase 3: Cross-chain edges}
\label{sec:cross-edges}

The remaining task is to assign edges between elements of
\emph{different} chains and between chain elements and skeleton
nodes they do not directly descend from.

\begin{conjecture}[Cross-chain consistency]
\label{conj:cross-chain}
The cross-chain edge assignment can be made composition-consistent,
$\forall$-safe, and compatible with the triangle types witnessed
in $\mathcal{G}$.
\end{conjecture}

This conjecture is equivalent to the Extension Solvability
Conjecture of~\cite{TrianglePaper}: if the triangle types from the
completion graph provide a rich enough vocabulary of three-way
relational patterns, the disjunctive constraint network for
cross-chain edges is solvable.

\begin{remark}[What the trichotomy achieves]
The key contribution of the trichotomy is to \emph{isolate} the
cross-chain problem as the only remaining gap.  Previous approaches
entangled the cross-edge problem with the infinite-chain problem
and the demand-satisfaction problem.  By showing that:
\begin{itemize}[nosep]
  \item Demand satisfaction for DR/PO is automatic (Case 1),
  \item Demand satisfaction for acyclic PP/PPI resolves with
    bounded expansion (Case 2),
  \item Infinite chains are algebraically rigid (Case 3),
\end{itemize}
the entire soundness question reduces to the single problem of
consistent cross-chain edge assignment.
\end{remark}


%% ============================================================
\section{Computational Evidence}
\label{sec:computation}

We present the exhaustive computational verification supporting
the trichotomy.

\subsection{Methodology}

For each configuration $(n, k)$ with $n$ nodes and $k$ types, we:
\begin{enumerate}[nosep]
  \item Enumerate all composition-consistent RCC5 networks on $n$ nodes.
  \item Enumerate all surjective type assignments with $k$ types.
  \item For each demand structure, simulate the tableau blocking scenario:
    create ``seed'' nodes (one per type) with successors for each demand,
    and test whether a blocked node's demands are satisfied.
  \item Test at depths 1, 2, and 3 for demand resolution.
\end{enumerate}

\subsection{Results}

\subsubsection{Self-absorption (16 entries)}
Exhaustive check of $R \in \comp(S, R)$ for all $R, S \in \NR^{-}$:
\textbf{15 pass, 1 failure} ($\PPI \notin \comp(\DR, \PPI)$).

\subsubsection{Demand satisfaction in arbitrary graphs}

\begin{center}
\renewcommand{\arraystretch}{1.2}
\begin{tabular}{lrrr}
\toprule
\textbf{Configuration} & \textbf{Tested} & \textbf{Failures} & \textbf{Rate} \\
\midrule
$n{=}3$, $k{=}2$, arbitrary & 246 & 204 & 82.9\% \\
$n{=}4$, $k{=}2$, arbitrary & 12,824 & 12,552 & 97.9\% \\
$n{=}5$, $k{=}2$, arbitrary & 1,232,970 & 1,230,800 & 99.8\% \\
\bottomrule
\end{tabular}
\end{center}

Arbitrary composition-consistent graphs have massive failure rates.
This is expected---they lack the structural constraints of tableau
completion graphs.

\subsubsection{Tableau simulation (CF-constrained)}

\begin{center}
\renewcommand{\arraystretch}{1.2}
\begin{tabular}{lrrr}
\toprule
\textbf{Configuration} & \textbf{Scenarios} & \textbf{Stuck L1} & \textbf{Rate} \\
\midrule
$k{=}2$, free blocker rel & 5,472 & 0 & 0\% \\
$k{=}2$, CF-forced rel & 18,688 & 318 & 1.7\% \\
$k{=}3$, CF-forced rel & 236,652 & 2,532 & 1.1\% \\
\bottomrule
\end{tabular}
\end{center}

With free blocker relation: \textbf{zero failures}.  With
CF-constrained relation: all failures involve $\DR$ blocker
relation and $\PPI$ demand (the unique self-absorption failure).

\subsubsection{Depth-2 resolution (excluding PP/PPI cycles)}

\begin{center}
\renewcommand{\arraystretch}{1.2}
\begin{tabular}{lrrrr}
\toprule
\textbf{Config} & \textbf{Stuck L1} & \textbf{Resolved L2} & \textbf{Remain} \\
\midrule
$k{=}2$, no cycle & 60 & 60 & \textbf{0} \\
$k{=}3$, no cycle & 1,194 & 1,170 & 24 \\
$k{=}3$, depth 3 & 24 & 24 & \textbf{0} \\
\bottomrule
\end{tabular}
\end{center}

For 2 types: \textbf{100\% resolution at depth 2}.  For 3 types:
the 24 remaining are PPI chains of length~2, all resolving at
depth~3.  \textbf{Zero failures when expansion depth matches chain
length}.


%% ============================================================
\section{Discussion}
\label{sec:discussion}

\subsection{What this paper establishes}

\begin{enumerate}[nosep]
  \item The complete algebraic analysis of self-absorption in
    RCC5: exactly one failure ($\PPI$ under $\DR$).
  \item The rigid/diluting classification: $\PP, \PPI$ rigid
    ($\comp(R,R) = \{R\}$); $\DR, \PO$ diluting
    ($\comp(R,R) = \NR^{-}$).
  \item The trichotomy: no PP/PPI $\to$ direct; acyclic PP/PPI
    $\to$ bounded expansion; cyclic PP/PPI $\to$ infinite chains.
  \item Exhaustive computational verification of all three cases.
  \item Reduction of the soundness problem to cross-chain edge
    assignment alone.
\end{enumerate}

\subsection{What remains open}

The cross-chain edge assignment
(Conjecture~\ref{conj:cross-chain}) is the single remaining
obstacle.  This is equivalent to the Extension Solvability
Conjecture of~\cite{TrianglePaper}, which is supported by
exhaustive computation (68,276 models, zero genuine failures)
but not formally proved.

\subsection{Relation to the two-tier quotient approach}

The PO-coherent fragment paper~\cite{TwoTierPaper} handles
PP-chains explicitly using kernel unfolding and period descriptors.
Its soundness proof avoids the Henkin construction by checking
path-consistency as an explicit precondition.  The present paper's
approach is broadly compatible: the two-tier quotient handles the
infinite-chain case (our Case~3), while our Cases~1--2 handle the
finite demand satisfaction.

A complete decidability proof might combine:
\begin{itemize}[nosep]
  \item The present paper's trichotomy for demand classification,
  \item The two-tier quotient's chain construction for Case~3,
  \item The triangle-type filtering of~\cite{TrianglePaper} for
    cross-chain edges.
\end{itemize}


%% ============================================================
\section{Conclusion}
\label{sec:conclusion}

We have presented a direct approach to the soundness problem for
the $\ALCIRCC{5}$ tableau, bypassing the Henkin construction.  The
key insight is that the RCC5 composition table has exactly one
self-absorption failure ($\PPI$ under $\DR$), and this failure
manifests only in PP/PPI chain demands.  All other demands are
automatically satisfied in any open completion graph.

The trichotomy decomposes the soundness problem into three cases
of decreasing difficulty:
\begin{enumerate}[nosep]
  \item DR/PO demands: immediate (zero failures).
  \item Acyclic PP/PPI chains: bounded expansion (computationally
    verified, depth $\leq$ number of types).
  \item Cyclic PP/PPI chains: infinite but rigid (standard chain
    construction).
\end{enumerate}

The only remaining gap is the cross-chain edge assignment---the
Extension Solvability Conjecture of~\cite{TrianglePaper}---which
is shared with all previous approaches.  The contribution of this
paper is to show that \emph{everything else works}: demand
satisfaction, chain construction, and the finite skeleton are all
resolved.  The decidability of $\ALCIRCC{5}$ rests entirely on
whether cross-chain edges can be assigned consistently.


%% ============================================================
\section*{Acknowledgments}

This research was prompted by Michael Wessel, who introduced the
$\ALCIRCC{}$ family.  The computational experiments were essential
for discovering the trichotomy and the unique self-absorption
failure---a fact that would be difficult to notice from the
composition table alone.


%% ============================================================
\bibliographystyle{plain}
\begin{thebibliography}{10}

\bibitem{CompanionPaper}
Claude.
\newblock On the decidability of {$\mathcal{ALCI}_{\mathit{RCC5}}$} and
  {$\mathcal{ALCI}_{\mathit{RCC8}}$}: Quasimodels meet the patchwork
  property.
\newblock Discussion paper, March 2026 (revised).
\newblock \emph{Note}: type-elimination algorithm retracted (unsound);
  tableau soundness unproven.  Tableau termination and completeness
  remain established.

\bibitem{TrianglePaper}
Claude.
\newblock Triangle types and the extension gap: Towards decidability of
  {$\mathcal{ALCI}_{\mathit{RCC5}}$}.
\newblock Discussion paper, April 2026.

\bibitem{TwoTierPaper}
Claude.
\newblock Two-tier quotient models for the PO-coherent fragment of
  {$\mathcal{ALCI}_{\mathit{RCC5}}$}.
\newblock Discussion paper, April 2026.

\bibitem{Renz1999}
J.~Renz.
\newblock Maximal tractable fragments of the {Region Connection Calculus}: A
  complete analysis.
\newblock In {\em Proc.\ IJCAI}, pages 448--455, 1999.

\bibitem{Wessel2003}
M.~Wessel.
\newblock Qualitative spatial reasoning with the
  {$\mathcal{ALCI}_\mathit{RCC}$} family---{F}irst results and unanswered
  questions.
\newblock Technical Report FBI-HH-M-324/03, University of Hamburg, 2002/2003.

\end{thebibliography}

\end{document}
